% !TeX root = ../papers/nc-p3.tex
\section{平面長さ四の一点を持つ十七支持の窓}
\label{sec:finite-planar-four-seventeen}
\begin{theorem}\label{thm:finite-seventeen-all}
全中央支持4の標準有限Frobenius窓で、四辺が有限で幾何学的支持が十七点以上であれば、
反対階数10で少なくとも片側対称性があり、反対階数11は生じない。
\end{theorem}
\begin{proof}[独立な鉛筆を持つ平面四重点]
前節で残るのは十六単純点と平面四重点 $d$ である。局所環は
$R=k[x,y]/(x^2,y^2)$ と書ける。実際 $\mathfrak m/\mathfrak m^2$ の積は
一次元のsocleへの非退化対称対合であり、二つの等方直線の基底を用いればよい。
$R$ での非零主イデアルの長さは単元なら4、非零一次部分があれば2、socleなら1である。
単純印 $p$ の五空間で $A(d),B(d)$ の剰余値が比例する印をまず数える側に回す。
残りの印では $d$ の環全体に独立な鉛筆があるので、非主イデアルによる差の除法は不要である。
$Z$ は前節までの $h(A)$ の零化イデアルの商とする。$d$ でのその因子は
長さ1,2の曲線状環、$R$、または空である。印以外の全単純点にも剰余二行列が比例すれば
外積は零で、その点は $Z$ に入らない。

有限な階数一切断で $J$ の外の十支持を受け取れるのは被約六点または単純四点と二重節点の型である。
後者では前節の計数で両超平面が節点を避ける。前者で超平面内の節点数が4以上なら
単独四空間の長さ4判定で対称性が得られる。2または3なら、混合割線を内部の同じ節点でまとめた
$\binom{k}{2}+k+\binom{6-k}{2}$ は9以下で十点に足りない。
従ってどちらの型でも超平面は節点間の割線や二重節点の接線を含まない。
$d$ での相対割線は定数か一つの長さ2節点を介するので、その超平面との交点である
$B$ の閉埋込みが長さ4を受け取れない。

曲線がある場合、母線を除くと全体を張る曲線の長さ4切断の議論はそのまま使える。
二次曲線と二つの被約孤立点の場合も $d$ では環全体で独立な鉛筆を使えるため、
定数孤立点の分離と固定核の直線による商外積の零化の証明が成立する。
前節の十支持の計数により、残るのは一つの埋込み接点と一つの孤立点がある場合である。
ここで補題\ref{lem:embedded-conic-contact-budget}の局所上界は任意のArtin基底 $T$ 上で使える。
単根が埋込み点に延びる場合、$z^2=0$ なら長さは $2\dim T/(z)$ 以下である。
二点因子の剰余が曲線の内側に接する場合、同補題の $D^2(A,B)=0$ より
\[
 \dim T=\dim T/(D^2)+\dim (D^2)
          \le 2\dim T/(D)+\dim T/(A,B).
\]
実際 $(D)/(D^2)$ は $T/(D)$ の商、$(D^2)$ は $T/(A,B)$ の商である。
前者の一次商 $T/(D)$ 上では鉛筆は曲線の張る平面に入り、後者 $T/(A,B)$ 上では固定接平面に入る。
こうして全長13の等号まで同じである。単独の非単純因子で固定核の直線に長さ2を使うので
$d$ の全長4を要する。内側に接する場合は $(A,B)$ が剰余で零だから接平面への長さも正となり、
等号は不可能である。残る単根の場合は $z^2=0$, $\dim R/(z)=2$ である。
従って $(z)$ は先の二つの等方方向の一つが定めるイデアルである。
socleは $(z)$ に入るので一次方向の選択で商は一意に決まる。
この場合の $p$ は核側でその長さ2部分の接線にあり、一つの側で高々二つの単純印しか得られない。
\end{proof}

\begin{proof}[母線に落ちる鉛筆とsocle]
行因子が一定の母線 $ku\otimes T_0$ を仮定する。列因子の場合は転置して同様である。
$\bar E$ と右評価を命題\ref{prop:curvilinear-marked-pencils}の通り取る。
$\dim\bar E=3$ のとき右評価の二次小行列のイデアルを $I$、
$\dim\bar E=2$ なら $I=(1)$ とする。$Z$ の各因子 $T$ で $l(u)I=0$ が成立すれば矛盾する。
実際 $(l(u))$ は $T/I$ の商なので $\dim T\le\dim T/(l(u))+\dim T/I$。
第一項の総和は印を通る平面で6以下である。第二項も同命題の評価の退化スキームにより6以下である。
長さ3の退化スキームの場合、各局所環を剰余体で組成列分解し、グラフへの引戻しの右完全性で
その逆像長を各方向の逆像長に組成長を掛けた和で抑えればよい。
これは曲線状でない $T$ でも成立する。曲線状因子では先の付値の不等式が $l(u)I=0$ を与える。

残る $T=R$ 上でこの積が非零になるには、鉛筆全体が剰余で $ku$ に値を取り、
$l(u)$ が非零socleに入る必要があることを示す。右評価には単元の軸があるので
射影した鉛筆の係数は一成分を消去した自由加群に入る。その階数は
$\dim\bar E-1\le2$ で、階数2の場合は残る一行の評価式の二係数が $I$ を生成する。
射影した二本の一方が剰余で非零ならこれを基底元にできる。
第二の射影からその成分を引くと、残る一成分 $b$ は $bI=0$ を満たす
（自由加群が一次元なら $b=0$）。$b$ が単元なら $I=0$。
さもなくば元の行列の結合は剰余で非零かつ $ku$ に値を持ち、
$b$ を法としても $ku$ に値を取るため $l(u)\in(b)$ である。

二本とも剰余射影が零なら、上の消去後の射影係数を二列の行列 $P$ として全成分が
$\mathfrak m$ にある。元の二行列の $u$ 成分は剰余で独立な二行ベクトルであり、
二つの座標を取れば可逆な二次行列 $S$ になる。共通左核の式をこの二座標で読むと
$l(u) S+KP=0$ となる。ここで $K$ は射影した行列に残りの左核成分を掛けて同じ二座標を読む写像である。
$l(u)$ に非零の一次部分があれば、その部分を消さない接ベクトルで微分した等式から
$K$ の剰余は階数2でなければならない。従って $K$ は二行二列可逆であり、
$P=-l(u)K^{-1}S$。残りの右評価行が $P$ を消すので $l(u)I=0$ である。
これで積の非零は主張のsocleの場合に限られる。

この例外で $u$ の定める平面は $d$ の位置1の接平面である。
実際 $p_1$ からの射影での式 $l(u)$ は元のこの平面式の単元倍であり、$d$ で一次部分が零となる。
$d$ で二つの偶位置を固定した階数1の行列を $u_d v_d^t$ とすると $u$ は $u_d$ に比例する。
行列の左右因子が両奇位置の接平面の式であることは、隣接する二点の縮約消失を微分すれば従う。
例えば $w(d_0,-,q_2,q_3)=0$ の $q=d$ での微分では $d_3,d_0$ が隣接するため
$q_2$ を微分した項が零で、$u_d v_d^t$ が位置3の接方向を消す。もう一側も同じである。
\end{proof}

\begin{proof}[四重点における二つの母線]
今 $d$ を印に取り直す。母線の例外印 $p$ では $A^d(p),B^d(p)$ はともに
$u_d$ を行因子とする階数1行列である。それぞれ $A^d(d)=B^d(d)=u_dv_d^t$ と異なる。
よって両四空間に $u_d$ を固定した母線が一本ずつある。
単一四空間の階数一切断はこの場合その一本の直線だけである。従って同じ側で
列因子固定型の例外印を併せて持つこともなく、全ての母線の例外印と $d$ は
一方の変数側で同じ直線に入る。そのような単純印は高々二つである。

この場合 $d$ における二つの四空間の和は、対称性がなければ六次元以上である。
四次元なら伸長型の印における共通切断の判定
（定理\ref{thm:finite-rank-nine-all-marks}の証明）を使う。
五次元と仮定し、$u_d$ による射影を取る。各単一四空間の射影には全非零元が
階数2の二次元鉛筆がある（定理\ref{thm:finite-marked-ruling-pencil}の冒頭の計算）。
$d$ を通るこの側の変数の平面 $J_d$ の外には十以上の単純点があり、
その各点で共通左右核を持つ二つの行列は独立である。$l(u_d)=0$ の点では、
階数2の射影鉛筆の非零元は左核を持たないから各行列は対応する母線上にある。
これで高々二点である。$l(u_d)\ne0$ の点では共通鉛筆が射影後も独立である。
射影空間が二次元なら共通右核がなく不可能である。三次元ならその射影平面の
階数一切断は有限である（正次元なら全非零元が階数2の射影直線と交わってしまう）。
従って共通右核の方向は長さ3の平面切断の高々三本の割線から得られ、
各方向は印以外に高々二単純点である。合計八点では足りない。

最後に対称性がないとし、各単純印で直交性から次元5以下の側を一つ選ぶ。
単純印の共通四空間なら既に矛盾する。各側で $d$ で剰余値が比例する印の集合を取ると、
その集合と $d$ は $d$ で取り直した両四空間の交叉に入る。
交叉が平面以下ならこれだけで単純印は高々五つである。実際 $d$ を通る任意の平面は
$R$ で少なくとも長さ2を切り、全長は7以下である。母線の例外があれば先の和の下界から
交叉は直線以下であり、比例する単純印は高々二つに落ちる。
よって母線例外のない側で選べる印は高々 $5+2=7$、
ある側でも高々 $2+2+2=6$ である。二側合わせても十六単純印を覆えない。
これで残る平面四重点の場合の対称性と反対階数11の排除が示された。
\end{proof}
